% bare_jrnl.tex
% V1.4b
% 2015/08/26
% by Michael Shell
% see http://www.michaelshell.org/
% for current contact information.                                                                 
%%%%%%%%%%%%%%%%%%%%%%%%%%%%%%%%%%%%%%%%%%%%%%%%%%%%%%%%%%%%%%%%%%%%

%% This is a skeleton file demonstrating the use of IEEEtran.cls
%% (requires IEEEtran.cls version 1.8b or later) with an IEEE
%% journal paper.
%%
%% Support sites:
%% http://www.michaelshell.org/tex/ieeetran/
%% http://www.ctan.org/pkg/ieeetran
%% and
%% http://www.ieee.org/


\documentclass[journal]{IEEEtran}

\usepackage{circuitikz}          % antes de babel
\usepackage[spanish, es-noshorthands]{babel}
\usepackage{amsmath}
\usepackage{float}
\usepackage{instructivo}
\graphicspath{{./}{./fig/}}
\usepackage[natbib=true,style=trad-unsrt,backend=biber]{biblatex}
\addbibresource{literatura.bib}

\hyphenation{op-tical net-works semi-conduc-tor}

\begin{document}
%
% paper title
% Titles are generally capitalized except for words such as a, an, and, as,
% at, but, by, for, in, nor, of, on, or, the, to and up, which are usually
% not capitalized unless they are the first or last word of the title.
% Linebreaks \\ can be used within to get better formatting as desired.
% Do not put math or special symbols in the title.
\title{Tarea 1: Prensa Hidráulica de Precisión}


\author{Matías~A.~Camacho,~\IEEEmembership{Estudiante,~TEC}
        ,~Juan~P.~Elizondo,~\IEEEmembership{Estudiante,~TEC}
        ~y~Sergio~Salazar,~\IEEEmembership{Estudiante,~TEC}
}


% The paper headers
\markboth{Análisis y Control de Sistemas Lineales, IS2026}%
{Análisis y Control de Sistemas Lineales}


% make the title area
\maketitle


\begin{abstract}
El resumen no debe de exceder de 150 palabras y debe establecer lo que fue hecho, como fue hecho, los resultados principales y su significado. No cite referencias en el resumen. No coloque ecuaciones, símbolos especiales o palabras claves.
\end{abstract}

% Se sugiere no más de cuatro palabras o frases cortas en orden alfabético, separadas por comas, que representen su reporte
\begin{IEEEkeywords}
Modelado, estados, Laplace, corriente, presión.
\end{IEEEkeywords}


%%%%%%%%%%%%%%%%%%%%%%%%%%%%%%%%%%%%%%%%%%%%%%%%%%%%%%%%%%%%%%%%%%%%%%%%%%%%%%%%%%%%%%%%%%%%%%
%%%%%%%%%%%%%%%%%%%%%%%%%%%%%%%%%%%%%%%%%%%%%%%%%%%%%%%%%%%%%%%%%%%%%%%%%%%%%%%%%%%%%%%%%%%%%%
\section{Introducción}

\IEEEPARstart{E}{n} esta sección se describe una introducción que le permita al lector identificar aspectos importantes para la comprensión de la lectura del documento, una visión relacionada al experimento y a los conocimientos científicos estudiados.


%%%%%%%%%%%%%%%%%%%%%%%%%%%%%%%%%%%%%%%%%%%%%%%%%%%%%%%%%%%%%%%%%%%%%%%%%%%%%%%%%%%%%%%%%%%%%%
%%%%%%%%%%%%%%%%%%%%%%%%%%%%%%%%%%%%%%%%%%%%%%%%%%%%%%%%%%%%%%%%%%%%%%%%%%%%%%%%%%%%%%%%%%%%%%
\section{Descripción dinámica}

...

\subsection{Circuito de control}

\subsection{Solenoide}
Un solenoide es un dispositivo electromagnético que convierte la energía eléctrica en movimiento mecánico. Para lograrlo, 
requiere de una bobina que genera un campo magnético cuando se hace fluir corriente por esta, atrayendo a un núclo ferromagnético
que varía su movimiento según la intensidad del campo magnético generado. 


\begin{figure}[H]
  \centering
  \begin{circuitikz}[american]
  \draw
    (0,0) to[V, v=$V_s$, invert]  (0,3)    % Fuente: lado izquierdo, sube
          to[R, l=$R$]            (3,3)    % Resistencia: hilo superior
          -- (5,3)                         % esquina superior derecha
          to[L, l=$L$]            (5,0)    % Bobina: lado derecho, baja
          -- (0,0);                        % hilo inferior

  % Flecha de corriente en el hilo inferior, bien alejada
  \draw[->, blue] (2.5,0) -- (1.0,0)
        node[midway, below] {$i(t)$};
\end{circuitikz} 
  \caption{Circuito equivalente de un solenoide.}
  \label{fig:circuito_solenoide}
\end{figure}

La figura anterior, muestra el circuito eléctrico equivalente para un solenoide como el que se va a modelar
en este trabajo. La resistencia $R$ representa la resistencia del alambre propio de la bobina, y la inductancia $L$
representa la inductancia de la bobina, la cual vista desde el núcleo ferromagnético, 
es una inductancia variable que depende de la posición del núcleo.

\setlength{\parskip}{1em}

A continuación, se muestran las ecuaciones que describen el comportarmiento 
eléctrico de este circuito \cite{Alexander2018}.

\begin{equation}
\begin{aligned}
    V_s &= i(t)R + L \frac{di}{dt} + i(t) \frac{dx}{dt} \frac{dL(x)}{dx} \\
        &= iR + L \dot{i} + i \dot{x} \frac{dL(x)}{dx}
\end{aligned}
\label{eq:eléctrica_solenoide}
\end{equation}

Esta primera ecuación relaciona la tensión aplicada en la entrada, con la corriente que fluye 
por el circuito. Si se observa, además de la aplicación de la ley de Ohm para el estado estacionario,
se consideran las variaciones en el tiempo de la corriente, así como el último término que toma en cuenta 
a la tensión "back-EMF" electromagnética, generada por el movimiento del émbolo.

\setlength{\parskip}{1em}

Así entonces, en el estado estacionario en el que no existe movimiento del émbolo, la ecuación se reduce
a únicamente el siguiente término. No obstante, para efectos del análisis dinámico del modelo y su linealización, 
se van a considerar los tres términos completos.

\begin{equation}
    V_s = i(t)R
\end{equation}

En donde además, se sabe que la resistencia depende del material propiamente, de la longitud $l$ y del 
área transversal $A$ del alambre de la bobina. 

\begin{equation}
    R = \rho \frac{l}{A}
\label{eq:resistencia_solenoide}
\end{equation}

Finalmente, la fuerza magnética generada por la bobina sobre el émbolo, se puede describir por medio
de la siguiente ecuación \cite{sadiku2001electromagnetics}.

\begin{equation}
    F_e = \frac{1}{2} i^2 \frac{dL(x)}{dx}
\label{eq:fuerza_magnetica_solenoide}
\end{equation}

En donde se ve claramente la dependencia con la corriente, así como la variación
de la inductancia respecto a la posición del émbolo \cite{mathworks_solenoid}. Este último término se puede reescribir,
de manera que la ecuación queda como.

\begin{equation}
    F_e = \frac{i^2}{2} \frac{-\beta}{(\alpha + \beta x)^2}
\label{eq:fuerza_magnetica_solenoide_reescrita}
\end{equation}


\subsection{Sistema hidráulico}
Los sistemas hidráulicos se basan en el principio de Pascal, el cual establece 
que un cambio de presión aplicado a un fluido confinado se transmite sin 
pérdidas a todo el fluido y a las paredes del recipiente que lo contiene 
\cite{PhysicsBook}. Para un fluido incompresible en equilibrio, la presión 
satisface:
\begin{equation}
    p = \frac{F_1}{A_1} = \frac{F_2}{A_2}
\label{eq:pascal}
\end{equation}
donde $F_1$ y $F_2$ son las fuerzas aplicadas sobre los pistones 1 y 2, de 
áreas $A_1$ y $A_2$ respectivamente. Esta relación permite que una fuerza 
modesta aplicada sobre el pistón de menor área se amplifique en el pistón de 
mayor área, lo cual constituye el principio de funcionamiento de la prensa.

\setlength{\parskip}{1em}

Adicionalmente, para un fluido incompresible, la conservación del volumen 
desplazado acopla los movimientos de ambos pistones:
\begin{equation}
    A_1 \, \Delta x_1 = A_2 \, \Delta x_2
\label{eq:conservacion_volumen}
\end{equation}
de modo que el desplazamiento del pistón 2 queda determinado por el 
desplazamiento del pistón 1, escalado por la razón de áreas.

\setlength{\parskip}{1em}

En esta aplicación, el émbolo del solenoide actúa como pistón 1 empujando el 
fluido, mientras que el pistón 2 entra en contacto con la carga a prensar. 
Puesto que la prensa está diseñada para actuar sobre un material físico 
(por ejemplo, una pieza a comprimir o estampar), dicha carga presenta una 
reacción mecánica ante el desplazamiento y la velocidad del pistón. Esta 
reacción se modela de forma genérica mediante un comportamiento viscoelástico
caracterizado por dos parámetros: la rigidez equivalente de la carga $K_c$, 
que representa la resistencia elástica del material a ser deformado (con 
unidades de N/m), y el coeficiente de amortiguamiento equivalente $B_c$, que 
representa las pérdidas viscosas o disipativas asociadas a la deformación 
(con unidades de N$\cdot$s/m). Ambos parámetros dependen del material 
específico sobre el cual opere la prensa y se determinan empíricamente o a 
partir de las propiedades mecánicas de dicho material. La fuerza de reacción 
que la carga ejerce sobre el pistón 2 queda entonces:
\begin{equation}
    F_2 = K_c \, x_2 + B_c \, \dot{x}_2
\label{eq:carga_hidraulica}
\end{equation}
donde $x_2$ y $\dot{x}_2$ son el desplazamiento y la velocidad del pistón 2, 
respectivamente.

Finalmente, se considera la contribución hidrostática debida a la diferencia 
de altura $h$ entre ambos pistones, la cual se establece por la geometría de 
la prensa y es, por tanto, constante:
\begin{equation}
    \Delta p_{h} = \rho \, g \, h
\label{eq:presion_hidrostatica}
\end{equation}
donde $\rho$ es la densidad del fluido y $g$ la aceleración de la gravedad.

%%%%%%%%%%%%%%%%%%%%%%%%%%%%%%%%%%%%%%%%%%%%%%%%%%%%%%%%%%%%%%%%%%%%%%%%%%%%%%%%%%%%%%%%%%%%%%
%%%%%%%%%%%%%%%%%%%%%%%%%%%%%%%%%%%%%%%%%%%%%%%%%%%%%%%%%%%%%%%%%%%%%%%%%%%%%%%%%%%%%%%%%%%%%%
\section{Análisis}

\subsection{Circuito de control}

\subsection{Solenoide}
Para este sub-sistema, se procede a realizar el análisis de la ecuación mecánica en base al 
siguiente diagrama de cuerpo libre. 

\begin{figure}[H]
  \centering
  \begin{circuitikz}[american]
    % Cuadrado
    \draw (0,0) rectangle (2,2);

    % Línea punteada vertical (x = 0)
    \draw[dashed] (1,-1) -- (1,3);

    % Etiqueta de la línea
    \node at (1.5,2.5) {$x = 0$};

    % Texto en el centro
    \node at (1,1) {$m$};

    % ===== Flechas lado izquierdo =====
    % Flecha superior (Fl)
    \draw[->] (-0.1,2) -- (-1.5,2)
        node[midway, above] {$F_l$};

    % Flecha media (Fd)
    \draw[->] (-0.1,1) -- (-1.5,1)
        node[midway, above] {$F_d$};

    % Flecha inferior (Fr)
    \draw[->] (-0.1,0) -- (-1.5,0)
        node[midway, above] {$F_r$};

    % ===== Flecha lado derecho =====
        \draw[->] (2.1,1) -- (3.5,1)
        node[midway, above] {$F_e$};
\end{circuitikz} 
  \caption{Diagrama de cuerpo libre del solenoide.}
  \label{fig:diagrama_cuerpolibre_solenoide}
\end{figure}

En donde $F_l$ corresponde a una posible perturbación del sistema, que para efectos prácticos se va a omitir posteriormente
en el análisis, $F_d$ representa la fuerza de rozamiento del émbolo, $F_r$ la fuerza del resorte que mantiene en un lugar
determinado (x = 0) al émbolo, y $F_e$ es la fuerza magnética descrita previamente en la sección anterior. La masa m, corresponde
a la masa del émbolo o pin propiamente.

Aplicando la sumatoria de fuerzas en el eje x, y sustituyendo las fuerzas por su representación dinámica ya descrita y conocida, 
se obtiene lo siguiente.

\begin{equation}
\begin{aligned}
        m \ddot{x} &= F_e - (F_d + F_r + F_l) \\
                   &= \frac{i^2}{2} \frac{dL}{dx} - (b \dot{x} + k x + F_l)
\end{aligned}
\label{eq:dinámica_solenoide}
\end{equation}

Como se puede observar en la ecuación \ref{eq:dinámica_solenoide}, la dinámica del sistema es no lineal, debido a la dependencia cuadrática 
de la corriente, así como a la dependencia de la posición del émbolo en el término de la fuerza magnética. Lo anterior, hace necesario
realizar una linealización para poder obtener un modelo de espacio de estados y función de transferencia, lo mismo ocurre para 
la ecuación \ref{eq:eléctrica_solenoide} debido a su tercer término.

\subsection{Sistema hidráulico}


Tomando el émbolo del solenoide como pistón 1, se tiene $\Delta x_1 = \Delta x$. 
Aplicando la conservación de volumen \eqref{eq:conservacion_volumen}, el 
desplazamiento del pistón 2 resulta:
\begin{equation}
    \Delta x_2 = \frac{A_1}{A_2} \, \Delta x
\label{eq:desplazamiento_pisto2}
\end{equation}
y derivando respecto al tiempo:
\begin{equation}
    \Delta \dot{x}_2 = \frac{A_1}{A_2} \, \Delta \dot{x}
\label{eq:velocidad_piston2}
\end{equation}

La fuerza de reacción de la carga sobre el pistón 2, según 
\eqref{eq:carga_hidraulica}, es entonces:
\begin{equation}
    F_2 = K_c \, \frac{A_1}{A_2} \, \Delta x 
        + B_c \, \frac{A_1}{A_2} \, \Delta \dot{x}
\label{eq:fuerza_piston2}
\end{equation}

Por el principio de Pascal \eqref{eq:pascal}, la presión en el fluido 
(despreciando momentáneamente la contribución hidrostática) es:
\begin{equation}
    p = \frac{F_2}{A_2} = \frac{A_1}{A_2^{\,2}} 
        \left( K_c \, \Delta x + B_c \, \Delta \dot{x} \right)
\label{eq:presion_pascal}
\end{equation}

Incorporando la presión hidrostática \eqref{eq:presion_hidrostatica}, la 
presión total de salida vista por la carga es:
\begin{equation}
    {\,p_{out} = \frac{A_1}{A_2^{\,2}} 
        \left( K_c \, \Delta x + B_c \, \Delta \dot{x} \right) 
        + \rho \, g \, h \,}
\label{eq:presion_salida}
\end{equation}

\setlength{\parskip}{1em}

De forma equivalente, la fuerza útil aplicada sobre la carga es:
\begin{equation}
    F_{out} = F_2 = \frac{A_1}{A_2} 
        \left( K_c \, \Delta x + B_c \, \Delta \dot{x} \right)
\label{eq:fuerza_salida}
\end{equation}

\setlength{\parskip}{1em}

De esta manera, el subsistema hidráulico queda caracterizado por una relación 
lineal entre el desplazamiento (y velocidad) del émbolo del solenoide y la 
presión (o fuerza) ejercida sobre la carga, manteniendo consistencia con las 
salidas del subsistema mecánico linealizado y facilitando la posterior 
obtención del espacio de estados y la función de transferencia del sistema 
completo.

%%%%%%%%%%%%%%%%%%%%%%%%%%%%%%%%%%%%%%%%%%%%%%%%%%%%%%%%%%%%%%%%%%%%%%%%%%%%%%%%%%%%%%%%%%%%%%
%%%%%%%%%%%%%%%%%%%%%%%%%%%%%%%%%%%%%%%%%%%%%%%%%%%%%%%%%%%%%%%%%%%%%%%%%%%%%%%%%%%%%%%%%%%%%%
\section{Desarrollo del espacio de estados}


%%%%%%%%%%%%%%%%%%%%%%%%%%%%%%%%%%%%%%%%%%%%%%%%%%%%%%%%%%%%%%%%%%%%%%%%%%%%%%%%%%%%%%%%%%%%%%
%%%%%%%%%%%%%%%%%%%%%%%%%%%%%%%%%%%%%%%%%%%%%%%%%%%%%%%%%%%%%%%%%%%%%%%%%%%%%%%%%%%%%%%%%%%%%%
\section{Desarrollo de la función de transferencia}


%%%%%%%%%%%%%%%%%%%%%%%%%%%%%%%%%%%%%%%%%%%%%%%%%%%%%%%%%%%%%%%%%%%%%%%%%%%%%%%%%%%%%%%%%%%%%%
%%%%%%%%%%%%%%%%%%%%%%%%%%%%%%%%%%%%%%%%%%%%%%%%%%%%%%%%%%%%%%%%%%%%%%%%%%%%%%%%%%%%%%%%%%%%%%
\section{Construcción del modelo simulado}




%%%%%%%%%%%%%%%%%%%%%%%%%%%%%%%%%%%%%%%%%%%%%%%%%%%%%%%%%%%%%%%%%%%%%%%%%%%%%%%%%%%%
%%%%%%%%%%%%%%%%%%%%%%%%%%%%%%%%%%%%%%%%%%%%%%%%%%%%%%%%%%%%%%%%%%%%%%%%%%%%%%%%%%%%
\section{Conclusiones}
Las conclusiones van aquí.

%%%%%%%%%%%%%%%%%%%%%%%%%%%%%%%%%%%%%%%%%%%%%%%%%%%%%%%%%%%%%%%%%%%%%%%%%%%%%%%%%%%%
%%%%%%%%%%%%%%%%%%%%%%%%%%%%%%%%%%%%%%%%%%%%%%%%%%%%%%%%%%%%%%%%%%%%%%%%%%%%%%%%%%%%
\appendices
\section{Pendiente}
El texto relacionado al apéndice va aquí.


\printbibliography

%%%%%%%%%%%%%%%%%%%%%%%%%%%%%%%%%%%%%%%%%%%%%%%%%%%%%%%%%%%%%%%%%%%%%%%%%%%%%%%%%%%%
%%%%%%%%%%%%%%%%%%%%%%%%%%%%%%%%%%%%%%%%%%%%%%%%%%%%%%%%%%%%%%%%%%%%%%%%%%%%%%%%%%%%
\section{Biografías de los autores}



 
\begin{IEEEbiographynophoto}{Matías A. Camacho}
        Estudiante del Instituto Tecnológico de Costa Rica (TEC) en la carrera de ingeniería electrónica.
        
        Correo: jeacamacho@estudiantec.cr
\end{IEEEbiographynophoto}

\begin{IEEEbiographynophoto}{Juan P. Elizondo}
        Realizó sus estudios de secundaria en el SNCCCR, sede UNA Región Brunca, y actualmente cursa la carrera de Ingeniería Electrónica en el Instituto Tecnológico de Costa Rica (TEC). 
        
        Fue estudiante de la Universidad de Costa Rica (UCR) durante el año 2022 y participó en programas de estudio en matemática en la Universidad Nacional (UNA) durante los años 2020 y 2021. 
        
        Cuenta con preparación y/o experiencia en áreas como:
        \begin{itemize}
            \item Arquitectura básica de redes, CISCO CCNA V7 (ITN), (2021).
            \item Fundamentos de ciberseguridad, CISCO Systems, (2022).
            \item Programa de tutorías estudiantiles, Tecnológico de Costa Rica, (2024).
            \item Diseño de circuitos digitales, TEC (2025).
        \end{itemize}
        
        Correo: juelizondo@estudiantec.cr
\end{IEEEbiographynophoto}



\vfill

\end{document}